% Compile with:
% lualatex -shell-escape test-babelbidi-minted.tex

\documentclass{article}

% babel is loaded internally by babelbidi.
\usepackage[
english,
main=persian
]{babelbidi}

\settextfont{HM XNiloofar}
\babelfont{tt}[persian]{Vazir Code}

\usepackage{minted}

\begin{document}



\section*{آزمون babelbidi و minted}

متن فارسی برای آزمون تعامل زبان راست‌چین با محیط‌های کد.

\subsection*{حالت پیش‌فرض}

\begin{minted}{text}
        Default: alpha = 1; beta = 2.
        {braces} [brackets] (parentheses)
    \end{minted}

\subsection*{\texttt{textdir=ltr}}

\begin{minted}[textdir=ltr]{text}
        LTR block: alpha = 1; beta = 2.
        {braces} [brackets] (parentheses)
    \end{minted}

\subsection*{\texttt{textdir=rtl}}

\begin{minted}[textdir=rtl]{text}
        RTL block: alpha = 1; beta = 2.
        {braces} [brackets] (parentheses)
    \end{minted}

\subsection*{\texttt{textlang=persian}}

\begin{minted}[textlang=persian]{text}
        این یک کد persian است.
    \end{minted}

\subsection*{\texttt{chardir}: جداکننده‌ها}

\begin{minted}[
        textdir=rtl,
        chardir={
            {\{}{l},
            {\}}{l},
            {[}{l},
            {]}{l},
            {(}{l},
            {)}{l}
        }
        ]{text}
        A1 Delimiters: B1
        A2 {curly braces} B2
        A3 [square brackets] B3
        A4 (parentheses) B4
    \end{minted}

\end{document}
